\section{Discussion}\label{sec:discussion}

The four analyses answer different documentary questions. The census and random assessment tested one-sample-per-participant criteria, the Workbench audit tested one factor-metadata endpoint, and the pre-verification tested sample classification. The results apply to their respective frames and do not estimate repository-wide prevalence.

\paragraph{Participant linkage.} The nine tranche studies with documented repeated sampling had no qualifying participant-to-sample mapping. Without it, an analyst cannot account for within-participant dependence. At the study level, counting accessions would count four ionisation-mode assays of one sample set (MTBLS1122, MTBLS1124, MTBLS1129, MTBLS1130) as four validation units. Two pre-verification deposits (MTBLS14358, MTBLS14376) shared a title and cohort description, but their documentation did not establish whether their samples overlapped.

\paragraph{Sample classification.} The census classified every row as quality control in 220 studies, including 183 where every sample-type value was \texttt{blank}. The 19 human pre-verification deposits came from that subset. Other fields identified study samples in 17 and participant specimens in 16. The label appears to be a template default. Omitting pre-filled values would reduce this ambiguity. The inconsistent assay-to-sample linkage in MTBLS14421 was detectable from the deposit but could not be resolved from it.

\paragraph{Evidence standards and limitations.} Unique labels alone would have qualified hundreds of census studies, although the random tranche included 9 studies with repeated sampling and 3 with cell lines. Screening should prespecify affirmative evidence and report failure categories. The documentary bound applies only to the 371 filtered candidates. Publication text was unavailable for 11 of 40 sampled studies. STRICT requires one sample per participant, so zero STRICT does not imply zero reconstructible studies. The Workbench result applies only to the audited endpoint in the selected studies. The author review supports the zero-STRICT finding but is not independent human validation. Component A followed exposure to a second model's answers, and shared negative Q4 assessments do not measure sensitivity for repeated sampling. Component B was a non-blind check by the original adjudicator. The frozen implementation limitations are reported in Section~\ref{sec:res-limits}. Missing documentation does not establish true sample dependence.

\paragraph{Recommendations.} Repositories should require a participant identifier or subject-to-sample mapping for each biological sample, flag deposits that share sample sets, and avoid pre-filled sample-type values. STROBE-MetEpi, currently under review, asks authors to report sampling, collection, validation populations, counts, timing and repository access, but does not require a machine-readable mapping from deposited samples to participants~\citep{kelly2026strobe}. A deposit-level mapping would complement those manuscript recommendations and support reuse of repeated measurements. Methodological studies such as this audit fall outside the guideline's scope.
